\section{What is implemented, and what remains}
\label{sec:status}

The requested destination is a solver suite that can supply credible nozzle
flow, performance, wall loading and sensitivities to the engine optimizer.
This requires an explicit operating domain and a shared physical state. A
successful MOC design calculation is one component of that suite.

\textbf{Completion count: zero of four at the requested integrated-MOC scope.}
Three areas have reusable numerical or modeling foundations. Shock fitting
has relevant local literature but no identified implemented fitted-shock
path. This count does not erase the functionality of the existing Q1D,
empirical-screening or JAX modules.

\begin{longtable}{@{}P{0.20\linewidth}P{0.36\linewidth}P{0.36\linewidth}@{}}
\toprule
\textbf{Requested capability} & \textbf{What exists} & \textbf{What is still required}\\\midrule
\endhead
Consistent variable-gas MOC &
Frozen-composition, variable-\(c_p\) Q1D expansion; property provenance,
bounded interpolation, enthalpy/entropy integration and sonic closure.
CEA chamber-property surfaces also exist. &
Consistent 2D characteristic equations, gas-dependent starting data,
node state and fluxes; separate equilibrium and finite-rate formulations.
Existing MOC remains calorically perfect.\\
Shock fitting &
Fold/coalescence detection and separation screens identify limitations.
They do not construct a shock. &
Unknown shock geometry, conservative jumps, downstream compatibility,
entropy/composition transport and treatment of subsonic regions.\\
Viscous/separation prediction &
Bartz-based thermal tools, a displacement screening estimate, four empirical
FSS-onset criteria, and an MDO separation margin. &
An integral or resolved nozzle boundary layer coupled to the core,
verified displacement/heat/friction effects, and a defined separated-flow
model. RSS, hysteresis and side loads are not supplied by the screens.\\
General differentiable forward-wall net &
NumPy bounded forward analysis; differentiable throat-arc kernel;
differentiable selected residuals and fixed-state partials. &
Full wall/axis/reflection/start-line/exit dependence, topology-valid
derivatives and total engine sensitivities.\\
\bottomrule
\end{longtable}

The implementation assessment is based on current source and call-path
inspection, with hashes in \code{implementation_inventory.json}. It is not
a new numerical qualification run. Previously completed MOC evidence is
used only where its recorded source hashes still match.

\subsection{What the recent MOC result establishes}
The documented constant-\(\gamma\), inviscid Nozzle-B model passes its
internal gates and a 32/48-to-48/64-point physical-coordinate refinement
comparison. The maximum angle difference is \(0.0228912^\circ\), minimum
field coverage \(0.950270\), and maximum Mach difference \(0.00080366\).
The accepted profile uses exact D/E attachment, explicit kernel capacity
and accurate seed controls. The complete pair took approximately
1,498 seconds; the added refined solve alone took approximately 1,213 seconds.
These are measured costs for that experiment, not representative timing
statistics for the whole solver.

That result demonstrates a usable bounded numerical configuration. It
does not establish variable-gas physics, a fitted shock, viscous interaction,
total wall-design derivatives, or a domain-wide default. The published
Rao example also uses different throat/start-line assumptions. Its
constant-gas variational conditions cannot simply be reused as a proof of
optimality for a reacting or viscous design \cite[PDF pp.~2--5]{rao1958}.

\subsection{Three different problems need separate interfaces}
\begin{enumerate}
\item \textbf{Forward analysis:} prescribe a physical wall, inlet state and
back pressure; determine flow, loads and performance.
\item \textbf{Inverse or variational design:} prescribe performance or
boundary conditions and determine a wall or control surface. The existing
Rao construction belongs here.
\item \textbf{Engine MDO:} vary chamber, geometry, cooling, feed and hardware
variables while closing the coupled engine state and its constraints.
\end{enumerate}
The optimizer needs the derivative of its actual analysis and state closure.
A gradient of an already optimized Rao control surface with respect to
independent solved nodal values is a different object.

\section{Current architecture and the default-solver problem}
\label{sec:architecture}
\subsection{The important live code paths}
\begin{longtable}{@{}P{0.34\linewidth}P{0.58\linewidth}@{}}
\toprule\textbf{Code location} & \textbf{Observed role}\\\midrule
\endhead
\code{raosim/rao_variational.py:solve_rao_bvp} &
Constant-gas control-surface BVP, backward reconstruction, independent
mesh gates and reliability reporting.\\
\code{raosim/moc_bounded.py:march_bounded_net} &
Host forward net with wall/axis handling, physical-strip diagnostics and
explicit incomplete/fold outcomes.\\
\code{raosim/frozen_flow.py:FrozenIdealGasTable} &
Fixed-composition thermally perfect property law and Q1D reference
expansion. Result scope explicitly excludes MOC/Rao characteristics.\\
\code{raosim/physics.py:boundary_layer_displacement} &
Screening estimate labeled repository heuristic, without an established
nozzle validity range; not an integral boundary-layer solve.\\
\code{raosim/separation.py:check_separation} &
FSS-onset screening from Q1D area-station states, or separate frozen Q1D
states; RSS remains unavailable.\\
\code{raosim/jax/moc_kernel.py} &
Padded throat-arc kernel; no general differentiable bounded-wall solver
or complete differentiable starting-line construction.\\
\code{raosim/jax/sensitivities.py} &
Fixed-state control-surface derivatives and explicit parameter partials;
total design and bell-wall dependencies remain deferred.\\
\code{raosim/mdo/nozzle.py:resolve_contour_provider} &
Only the analytic TOP/B\'ezier provider is registered. The typed provider
outcome supports explicit availability/infeasibility.\\
\code{raosim/mdo/contour_validation.py} &
Postsolve comparison against a variational result; it does not replace
the optimizer's inner-loop geometry or thrust law.\\
\bottomrule
\end{longtable}

\subsection{Geometry substitution alone does not change engine aerodynamics}
The current \code{engine_residual} in \code{raosim/mdo/engine.py}
evaluates chamber surfaces and closes the scaled throat radius and mass
flow using
\begin{align}
R_1 &=
\frac{F_{\mathrm{req}}-\eta_{CF}C_{F,\mathrm{1D}}P_c\pi \Rt^2}
     {F_{\mathrm{req}}},\\
R_2 &=
\frac{\dot m-P_c\pi\Rt^2/(\eta_{c^*}c^*_{\mathrm{ideal}})}
     {\dot m_{\mathrm{ref}}}.
\end{align}
These are a transcription of the present software closure, not a proposed
general-gas formula. In \code{solve_engine}, ideal \(C_F\), exit Mach,
exit pressure and separation margin are computed before
\code{evaluate_contour} supplies the cooling/mass grid.

Consequently, registering an MOC geometry provider alone can change wall
geometry, cooling and mass while thrust, throat size and their sensitivities
continue to use the older Q1D performance model. A default transition must
update the coupled aerodynamic state contract and engine residual.

\begin{figure}[htbp]
\centering
\begin{tikzpicture}[
box/.style={draw=teal,rounded corners=2pt,fill=pale,text width=6.4cm,
align=center,font=\small,inner sep=6pt},
arr/.style={-{Latex[length=2mm]},thick,draw=muted},
node distance=6mm]
\node[box] (a) {Current design variables and chamber surfaces};
\node[box,below=of a] (b) {Q1D \(C_F,c^*\) close \(\Rt,\dot m\)\\and define nozzle performance};
\node[box,below=of b] (c) {TOP/B\'ezier provider supplies geometry};
\node[box,below=of c] (d) {Cooling, feed, pumps, mass\\and engine metrics};
\node[box,right=10mm of a] (e) {Proposed design variables, gas mode\\and declared applicability};
\node[box,below=of e] (f) {Coupled flow + sizing residuals\\with consistent mass/energy/thrust};
\node[box,below=of f] (g) {One flow/performance/wall-state result\\with quality and derivative status};
\node[box,below=of g] (h) {Cooling, feed, pumps, mass\\and coupled objective/constraint totals};
\draw[arr] (a)--(b); \draw[arr] (b)--(c); \draw[arr] (c)--(d);
\draw[arr] (e)--(f); \draw[arr] (f)--(g); \draw[arr] (g)--(h);
\end{tikzpicture}
\caption{Current dependency order and proposed integration target.
The right-hand structure is a design proposal; it is not the present MDO implementation.}
\end{figure}

\subsection{Three existing \(c^*\) conventions must be reconciled}
\begin{longtable}{@{}P{0.30\linewidth}P{0.32\linewidth}P{0.30\linewidth}@{}}
\toprule\textbf{Path} & \textbf{Current computation} & \textbf{Integration implication}\\\midrule
\endhead
Host \code{cea.cea_propellant} &
Obtains RocketCEA \code{get_Cstar}. &
Chemical mode, database and station conventions must be preserved.\\
Host \code{frozen_flow} &
Computes \(c^*=p_0/G_*\) from its integrated variable-\(c_p\) sonic state. &
A physically distinct fixed-composition Q1D model.\\
MDO \code{ChamberSurfaces.c_star_ideal} &
Uses the calorically perfect analytic law evaluated with interpolated
\(\gamma,T_c,R\). The sampled CEA fields are these three quantities,
not \(c^*\) or \(G_*\). &
Sampling chamber properties does not automatically import the full
CEA or variable-property sonic closure into MDO.\\
\bottomrule
\end{longtable}

\subsection{A proposed shared flow-result contract}
\proposal Introduce one result consumed by sizing, cooling, structural
screens and reporting. It should include the physical wall and effective
core geometry; inlet and local gas-mode identifiers; mass flux; thrust and
its control-volume basis; wall-edge \(p,T,\rho,V,M\), composition and
transport properties; shock/boundary-layer state where applicable; and
reference units and station conventions.

Availability, physical-domain status, conservation residuals, resolution
evidence and derivative validity must be separate fields. A finite,
fixed-shape array used to carry an infeasible candidate through JAX is
not a valid flow field. The present typed contour-provider behavior is
useful infrastructure for this distinction.

No explicit fallback should be hidden inside a higher-fidelity result.
Lower-fidelity analysis can be selected or used to initialize another
solver, with the executed model recorded in the result and optimizer
history.
